\chapter{Storage and Identity}
\label{ch:pebble-storage}

A compiler can preserve every arithmetic identity and still fail because an edit invalidates a pointer into a growing vector.
Memory policy is therefore part of the IR contract.
The initial PebbleIR policy is unusually plain: one caller-owned object contains all function storage, and IDs refer to slots rather than addresses.

\section{A bounded arena without allocation}
\label{sec:pebble-arena}
The function's fixed arrays act as a \term{region}: all value records share the function's lifetime.
Creating a value consumes a slot; deleting it leaves a tombstone; destroying the function releases the region at once.
There is no per-node allocation and no individual free.
The public API never returns a pointer to a value record, so clients need not keep interior pointers stable across a whole-function copy.

A caller can put a \code{Pir} on the stack, in static storage, or in a separately allocated owner object.
Those choices have different resource limits.
The teaching tests use stack storage, including temporary transaction copies; an embedded compiler with a small stack should put those workspaces in caller-owned heap storage instead.
The absence of allocation inside the API does not imply zero memory cost.

\section{Publication has one commit point}
\label{sec:pebble-publish}
Listing~\ref{lst:pebble-allocation} is the actual allocation path.
It checks all local failure conditions before changing either the count or the revision.
Once those checks pass, storing a bounded record cannot fail in this representation.
The builder checks only local construction conditions; the supplied opcode and operands still need full verification.

\begin{listing}[htbp]
\inputminted[firstline=22,lastline=33]{c}{example-code/pebbleir/pebbleir.c}
\caption{A failed value allocation leaves the function unchanged.}
\label{lst:pebble-allocation}
\end{listing}

The \term{revision} increases on a committed edit.
Derived facts carry both the owning object pointer and the revision they describe.
An observer must compare both before reusing a cached answer.
The counter is checked against \code{ULONG\_MAX}; wrapping it could make a very old result appear current.

Initialization, successful parsing, and successful source compilation begin a new logical function lifetime.
They can reset the revision, so every prior fact and ID association for the destination must be discarded even if its C address is reused.
A long-lived service should add an owner-generation token or globally unique function identity if such lifetime transitions cannot be enforced by its calling convention.

\begin{algorithm}[htbp]
\caption{Replacing a complete function as a transaction.}
\label{alg:pebble-transaction}
\begin{pseudocode}
\Require destination $F$, candidate-producing operation $P$
\Ensure verified replacement or unchanged destination
\State construct candidate $T$ in separate storage
\If{construction fails or verification rejects $T$}
\State discard $T$ and report the failure
\Else
\State invalidate facts belonging to the old lifetime
\State publish $T$ into $F$
\EndIf
\end{pseudocode}
\end{algorithm}

Algorithm~\ref{alg:pebble-transaction} is used by parsing and source compilation.
Optimization instead copies an existing function, increments its revision, verifies the candidate, and commits it within the same logical lifetime.
The distinction matters to cache owners even though both implementations end with a C structure assignment.

\section{Worked capacity boundary}
\label{sec:pebble-full}
With 127 occupied slots, allocating a constant returns ID 128 and increments the revision once.
The next allocation returns zero and changes neither the count nor any slot.
The test suite fills the table and compares the object before and after the failed request.
This is an interface property, not merely a check that a sanitizer sees no overflow.

For growing programs, retain the IDs while replacing the fixed arrays with growable storage.
Check multiplication before computing allocation sizes, define behavior on allocation failure, and publish a new buffer only after successful allocation and copying.
If deleted IDs are reused, attach generation numbers so an old handle cannot accidentally name a new instruction in the same slot.
These are explicit extensions; the supplied implementation never reuses a slot within a function lifetime.

\section{Exercises}
\label{sec:pebble-five-exercises}
\begin{enumerate}
\item \label{ex:pebble-revision} A fact has the correct owner pointer but an old revision. Can it be used if the last edit merely changed a constant? Give a client for which the answer is unsafe.
\item \label{ex:pebble-generation} Describe a stale-handle failure when slot 12 is deleted and reused. What additional field distinguishes the old and new references?
\end{enumerate}

\section{Further reading}
\label{sec:pebble-five-reading}
The C11 object-lifetime and allocation rules delimit what the host implementation may retain across storage replacement~\cite[Sections~6.2.4 and~7.22.3]{wg14-2011-c11}.
